\chapter{Discussion and Future Work}\label{chapter:discussion}

The techniques reviewed in Chapter~\ref{chapter:techniques} address capacity pressure, access bandwidth pressure, and state preparation overhead through different design choices. We discuss their interactions and two directions motivated by agentic applications, in which an LLM makes repeated model calls and uses external tools to pursue a goal: longer contexts within individual requests and longer task lifetimes across requests.

\section{System-wide Trade-offs and Co-design}\label{sec:discussion:codesign}

Optimizations in one part of the infrastructure change the benefits of others. Quantization reduces storage and transfer volume but adds conversion work; offloading expands capacity but can increase preparation delay. Sparse access saves bandwidth only when the storage layout allows selected entries to be delivered efficiently. Likewise, GQA and MLA reduce the size of KV relative to hidden states, potentially weakening the capacity advantage of restoration methods such as HCache~\cite{kv_gao2025hcache}. Benefits therefore depend on the combined design and cannot simply be multiplied across techniques.

Co-design should connect model representations, physical access, and execution management. HiSparse illustrates how model-level sharing of sparse selections can expose future accesses and enable earlier transfers~\cite{kv_xie2026hisparse}. More broadly, reducing cache size is useful when it improves execution under the relevant resource constraints. Conversion, selection, and restoration costs may shift the dominant bottleneck. The three challenges provide a common basis for examining these interactions, including the quality costs introduced when several approximate methods are combined.

\section{Ultra-Long Contexts}\label{sec:discussion:long-context}

Complex agentic tasks may require a single request to incorporate extensive documents, code, tool outputs, and reasoning history. Supporting such contexts requires both sufficient capacity and affordable access. With dense attention, every decoding step reads historical KV, so moving more history into a larger storage tier can leave bandwidth as the limiting factor. Future infrastructure should jointly develop compact representations and selective access that control the growth of data movement while preserving long-range dependencies. Native Sparse Attention provides an example of designing sparsity together with training and hardware-efficient execution~\cite{kv_yuan2025nsa}.

An additional challenge is that historical information can become useful much later. An agent may revisit early evidence after obtaining new tool results; low attention at the present step does not establish that this evidence can be discarded. Future designs should distinguish information that need not be accessed now from information that can be removed permanently, and preserve a means of recovering useful history. Task progress and the role of different context segments may help guide these choices. Long contexts also make preparation expensive: incremental construction and reuse should amortize this cost across calls, including when materials are added, reordered, or revised. The objective is efficient and sustained use of long histories, with capacity, access cost, preparation time, and task quality considered together.

\section{Long-lived Agentic Tasks}\label{sec:discussion:long-lived}

Agents may execute for hours or days, alternating model computation with periods of waiting for external tools. Task lifetime therefore extends beyond an individual inference request. During a wait, the context remains valuable although the model is temporarily inactive. Keeping its KV on the GPU consumes capacity that other tasks could use; discarding it makes the next call more expensive to prepare. Future infrastructure should separate the lifetime of retained state from its residence in scarce GPU memory, supporting inexpensive preservation and timely resumption. TokenCake already explores releasing GPU cache capacity during tool calls and preparing state for subsequent execution~\cite{kv_bian2025tokencake}.

Long-lived tasks also require continuity as their histories evolve or execution is interrupted. Context summarization and revision can invalidate previously computed KV, while failures can make extensive reconstruction necessary. Infrastructure should support reuse as the task changes and recovery at a cost appropriate to its accumulated work. GhostServe illustrates protecting KV for recovery after device failures~\cite{kv_jayakody2026ghostserve}. KV is nevertheless recomputable model state: resuming an agent also requires its recorded tool results and other task state. Their coordination is essential to preserving the intended context across interruptions. Longer contexts and longer lifetimes reinforce each other, making efficient preservation, access, and restoration central to infrastructure for complete agentic tasks.
